\section{基于集成的分类}
\label{sec:classification-ensemble}

集成可以组合硬标签、类别得分或概率，整体输出决定分类任务的使用方式。例如GBDT用回归树拟合类别得分的修正，再由累积得分形成类别决策；成员的数值输出与整体分类目标相容。下面比较各方法的输出、阈值和概率检验，成员组合的基本机制见第\ref{sec:classic-overview-ensemble}节，完整训练过程见第\ref{chap:ensemble-learning}章。

\begin{table*}[htbp]
  \centering
  \small
  \begin{tabularx}{\linewidth}{@{}>{\RaggedRight\arraybackslash}p{22mm}>{\RaggedRight\arraybackslash}X>{\RaggedRight\arraybackslash}X@{}}
    \toprule
    方法 & 分类输出 & 任务中的条件与代价 \\
    \midrule
    Bagging & 硬标签投票，或平均成员得分／概率后决策 & 概率平均不自动校准；预测成本随所用成员增加 \\
    随机森林 & 树投票或叶类别比例平均 & 树深、特征抽样与树数需选择；袋外诊断不替代测试 \\
    AdaBoost & 对弱分类器加权投票或输出累计得分 & 加权分类问题上的成员表现与噪声影响轮次选择 \\
    GBDT & 累加类别得分，再按阈值或Softmax决策 & 分类损失决定修正；树深、学习率与早停共同影响结果 \\
    Stacking & 元学习器组合标签、得分或概率 & 组合层用留出或折外基预测训练；校准仍需独立检验 \\
    \bottomrule
  \end{tabularx}
  \caption{集成方法的分类输出与使用条件。此处按分类输出与训练条件比较；成员训练机制见概述与集成学习章。}
  \label{tab:classification-ensemble-map}
\end{table*}

集成的投票和得分需要按整体预测评价。表\ref{tab:classification-evaluation-protocol}给出标签错误、排序、概率与误判代价的对应指标，资源成本计入实际使用的全部成员。组合后的概率质量和解释性由整体模型检验。对于奇数个二分类成员的等权多数投票，若在目标分布几乎所有输入与标签上，成员错误在给定输入与标签后关于训练随机性相互独立，且每个成员的条件错误率都至少比$1/2$小同一个正数，则平均训练随机性后的投票风险具有随成员数指数下降的上界。二分类AdaBoost的训练误差上界下降则需各轮保持带权优势，泛化还需投票间隔或基类容量控制（参考\BookSectionRef{sec:boosting}{卷3《基础、理论与可信性》第5章“二分类学习的可学习性”第5.6节“集成学习”}）。
